\chapter{Use Cases}
\ChapterLabel{UseCases}
\section{Chapter Responsibility}
\label{Section:UseCasesResponsibility}
Describe how actors interact with the system to achieve their goals.

\noindent\textbf{Expected result:} Complete interaction scenarios tied to requirements and tests.

\section{Revision 0.2 Scope Notice}
\label{Section:DeferredUseCases}
\begin{fsaNote}[Project use cases --- draft for review]
This chapter defines draft use cases linked to the project stories,
requirement candidates, and planned tests. Joel Murphy and Jason Lefler own
the open decisions and review the proposals.
See \hyperref[Table:ProjectGapReview]{the reviewed gap list} for actual
decisions and unresolved items.
\end{fsaNote}

\section{How to Use This Chapter}
Describe goal-oriented interactions at the system boundary. Give every use case
a stable ID, its own table, and at least one reciprocal requirement mapping.
Include preconditions, trigger, main flow, alternate or exception flows, and
postconditions. Maintain the canonical relationships among stories,
requirements, use cases, and tests in \ChapterRef{Traceability}.

\section[High-Level Use Case Diagram]{\hspace{0.5em}High-Level Use Case Diagram}
\label{Section:ProjectUseCaseDiagram}
Figure~\ref{Figure:UseCaseDiagram} shows the sensor-interface subsystem and
its draft actor goals. Each use-case ID in the diagram links to its
use-case table; each table links back to the diagram. Actor associations
show participation rather than signal flow. The subsystem boundary and
external dependencies remain defined in
\hyperref[Section:ProjectSystemBoundary]{System Context}.

\begin{figure}[htbp]
\centering
% Project UML use-case diagram. Each use-case ID links to its table row.
\begin{tikzpicture}[
  external actor/.style={draw=fsaGray,rounded corners=2pt,fill=fsaGray!8,
    align=center,text width=2.5cm,minimum height=1.15cm,font=\small},
  project use case/.style={ellipse,draw=fsaBlue,fill=fsaLightBlue,
    align=center,text width=3.8cm,minimum height=1.25cm,font=\small},
  association/.style={draw=fsaGray,line width=0.7pt}
]
  % Reserve a separate header band above the use-case ovals.
  \draw[fsaBlue,rounded corners=3pt] (-2.95,-2.0) rectangle (2.95,3.2);
  \node[align=center,font=\small\bfseries,text width=4.8cm] at (0,2.45)
    {Sensor-interface subsystem\\FTDI245R acquisition\\pipeline code base};
  \node[project use case] (acquire) at (0,0.65)
    {\UseCaseRef{UC-001}\\Acquire sensor data};
  \node[project use case] (outcomes) at (0,-1.0)
    {\UseCaseRef{UC-002}\\Handle acquisition\\read outcomes};
  \node[external actor] (engineer) at (-4.5,-0.15)
    {\textless\textless actor\textgreater\textgreater\\Embedded Engineer};
  \node[external actor] (consumer) at (4.6,0.65)
    {\textless\textless actor\textgreater\textgreater\\Downstream processing\\consumer};
  \draw[association] (engineer.east) -- (-3.05,-0.15) |- (acquire.west);
  \draw[association] (engineer.east) -- (-3.05,-0.15) |- (outcomes.west);
  \draw[association] (acquire.east) -- (consumer.west);
\end{tikzpicture}

\caption{Project UML use-case diagram linked to the draft use-case tables.}
\label{Figure:UseCaseDiagram}
\end{figure}
\index{use case diagram}

These drafts apply \hyperref[Skill:AIS-07]{AIS-07} and its
\hyperref[Section:AIUseCaseContext]{Use-Case Context}. They elaborate
\UserStoryRef{US-001}, \UserStoryRef{US-002}, \UserStoryRef{US-004}, and
\UserStoryRef{US-005}. The approved decisions in \hyperref[Section:ProjectAcquisitionConditions]{Requirements}
are reflected below. Remaining interface, timing and recovery details are TBD
(decision co-owners: Joel Murphy and Jason Lefler).

\section{Acquire Sensor Data}
\label{Section:ProjectAcquireUseCase}
\begin{tabularx}{\linewidth}{L{3.5cm}Y}
\toprule
\textbf{Identifier} & \UseCaseID{UC-001}\phantomsection\label{Table:ProjectUC001}; revision 0.2; Draft for review.\\
\textbf{Mapped story} & \UserStoryRef{US-001}.\\
\textbf{Mapped requirements} & \RequirementRef{REQ-F-008}; \RequirementRef{REQ-F-009}; \RequirementRef{REQ-F-010}; \RequirementRef{REQ-I-002}; \RequirementRef{REQ-D-002} (candidates; Open). Decision co-owners: Joel Murphy and Jason Lefler.\\
\textbf{Goal} & Receive data supplied through the FTDI245R interface and make it available to downstream processing unchanged and in order.\\
\textbf{Primary actor} & Embedded Engineer. Supporting external actor: downstream processing consumer.\\
\textbf{Preconditions} & Supported host with required driver/library installed and configuration completed; proposed FTD2XX. Input is one unsigned 8-bit byte per ADC0804 conversion, with no headers/framing. Concrete hardware, OS and driver/library versions, rate and consumer API remain TBD (decision co-owners: Joel Murphy and Jason Lefler).\\
\textbf{Trigger} & First connection of supported FT245R hardware after prerequisites are met; detection and acquisition start automatically without a code change.\\
\textbf{Main flow} & 1. The pipeline starts acquisition as described in \UseCaseRef{UC-003}.\par 2. The source supplies constant-rate raw sample bytes; USB delivery may be batched.\par 3. The pipeline makes positive-length reads and evaluates status/count.\par 4. It delivers successfully read raw bytes unchanged and in order across batches; delivery occurs when the consumer accepts bytes through its interface.\par 5. Acquisition continues until the user stops it.\\
\textbf{Alternate flow} & A1. Partial/empty timeouts follow \UseCaseRef{UC-002}; acquisition continues.\par A2. Known-sequence acceptance uses both a software-generated source and physical FT245R hardware. Selected Linux, Windows and macOS configurations are planned; execution is pending.\\
\textbf{Exception flow} & E1. A read failure or interruption follows \UseCaseRef{UC-002}.\par E2. Consumer unavailability/failure requires a policy; retry, discard, persistence, and reporting remain TBD (decision co-owners: Joel Murphy and Jason Lefler).\\
\textbf{Outcome} & Normal: the known source sequence is accepted unchanged and in order under the agreed conditions. After interruption, completeness is limited by the approved preservation boundary.\\
\textbf{Postconditions} & Only bytes successfully read into host memory are delivered. User-stop/drain details and consumer-error policy remain TBD (decision co-owners: Joel Murphy and Jason Lefler).\\
\textbf{Mapped tests} & \TestCaseRef{TC-010}; \TestCaseRef{TC-014}; \TestCaseRef{TC-016} (planned; not executed). Decision co-owners: Joel Murphy and Jason Lefler.\\
\textbf{Diagram / validation} & \hyperref[Figure:UseCaseDiagram]{UML diagram}; \ValidationRef{VAL-004} (known-sequence acquisition).\\
\textbf{Source / decisions} & TN\_167 Sections 2.1--2.2.1, 2.4 \cite{FTDI_TN167_v1_0}; D2XX Guide v1.3 Section 3.10 \cite{FTDI_D2XX_v1_3}. Constant sample rate value, acceptance-test duration, buffering, concrete consumer API and host/library versions remain TBD (decision co-owners: Joel Murphy and Jason Lefler).\\
\bottomrule
\end{tabularx}

\section{Handle Acquisition Read Outcomes}
\label{Section:ProjectReadOutcomesUseCase}
\begin{tabularx}{\linewidth}{L{3.5cm}Y}
\toprule
\textbf{Identifier} & \UseCaseID{UC-002}\phantomsection\label{Table:ProjectUC002}; revision 0.2; Draft for review.\\
\textbf{Mapped stories} & \UserStoryRef{US-002}; \UserStoryRef{US-012} (failed-read exception).\\
\textbf{Mapped requirements} & \RequirementRef{REQ-F-004}; \RequirementRef{REQ-F-005}; \RequirementRef{REQ-F-006}; \RequirementRef{REQ-F-007}; \RequirementRef{REQ-F-011}; \RequirementRef{REQ-F-012}; \RequirementRef{REQ-F-013}; \RequirementRef{REQ-Q-002}; \RequirementRef{REQ-D-002} (candidates; Open). Decision co-owners: Joel Murphy and Jason Lefler.\\
\textbf{Goal} & Distinguish completed, partial, empty-timeout, and failed reads so that interruptions do not crash or corrupt the pipeline or introduce invalid acquired data.\\
\textbf{Primary actor} & Embedded Engineer. Read outcomes arise at the controlled device, USB-stack, and host-platform boundaries.\\
\textbf{Preconditions} & An acquisition read is initiated through \UseCaseRef{UC-001}. Each request is at least one byte. Exact positive block size and timeout duration remain TBD (decision co-owners: Joel Murphy and Jason Lefler).\\
\textbf{Trigger} & An acquisition read is attempted and its completion, timeout, or error must be handled.\\
\textbf{Main flow} & 1. Configure the timeout through \texttt{FT\_SetTimeouts} and check status.\par 2. Attempt \texttt{FT\_Read} for at least one byte.\par 3. Recognize completion exactly when status is \texttt{FT\_OK} and actual count equals requested count.\par 4. Deliver successfully returned bytes unchanged and in order.\\
\textbf{Alternate flow} & A1. \texttt{FT\_OK} with a positive count below the request is a partial timeout: deliver those actual bytes in order, then continue acquisition.\par A2. \texttt{FT\_OK} with zero bytes is an empty timeout: deliver no bytes and continue acquisition.\\
\textbf{Exception flow} & E1. Timeout-configuration failure is reported; subsequent response is TBD (decision co-owners: Joel Murphy and Jason Lefler).\par E2. A failed read is distinct from an empty timeout; error status alone is not proof of disconnection.\par E3. Confirmed disconnection follows \UseCaseRef{UC-004}.\\
\textbf{Outcome} & Read outcomes are distinguished; timeouts continue acquisition. Disconnection preserves host-read bytes without crash, is reported as an interruption, and is followed by automatic resumption on reconnection.\\
\textbf{Postconditions} & Downstream bytes preserve value, count and order across batches, timeouts and disconnection within the host-read guarantee. Reporting mechanism, reconnect timing/identity, stale device bytes and consumer-failure handling remain TBD (decision co-owners: Joel Murphy and Jason Lefler).\\
\textbf{Mapped tests} & \TestCaseRef{TC-010}; \TestCaseRef{TC-011}; \TestCaseRef{TC-012}; \TestCaseRef{TC-013}; \TestCaseRef{TC-015} (planned; not executed). Decision co-owners: Joel Murphy and Jason Lefler.\\
\textbf{Diagram / validation} & \hyperref[Figure:UseCaseDiagram]{UML diagram}; \ValidationRef{VAL-005} (read outcomes and disconnection); \ValidationRef{VAL-014} (failed read versus empty timeout).\\
\textbf{Source / decisions} & D2XX Guide v1.3 Sections 3.10 and 3.15 \cite{FTDI_D2XX_v1_3}; owner-approved \hyperref[Section:ProjectAcquisitionConditions]{acquisition conditions} and operational definitions. Timeout duration, reporting interface, quantitative reliability and platform/library versions remain TBD (decision co-owners: Joel Murphy and Jason Lefler).\\
\bottomrule
\end{tabularx}

\section{Start Acquisition on First Connection}
\begin{tabularx}{\linewidth}{L{3.5cm}Y}
\toprule
\textbf{Identifier} & \UseCaseID{UC-003}; Draft for review.\\
\textbf{Mapped story} & \UserStoryRef{US-004}.\\
\textbf{Mapped requirements} & \RequirementRef{REQ-F-009}; \RequirementRef{REQ-F-010}.\\
\textbf{Goal} & Begin capturing samples by connecting supported hardware to a configured host.\\
\textbf{Primary actor} & Embedded Engineer.\\
\textbf{Preconditions} & The host is supported, the required driver/library is installed, and configuration is complete. See \hyperref[Section:ProjectStartupPrerequisites]{startup prerequisites}.\\
\textbf{Trigger} & The Embedded Engineer connects supported FT245R hardware for the first time in the scenario.\\
\textbf{Main flow} & 1. The Embedded Engineer connects the hardware.\par 2. The pipeline detects it without a code change.\par 3. The pipeline starts acquisition automatically and proceeds to \UseCaseRef{UC-001}.\\
\textbf{Alternate flow} & If the connection follows an interruption during acquisition, follow \UseCaseRef{UC-004}.\\
\textbf{Exception flow} & If a prerequisite is unmet or the hardware is unsupported, the automatic-start guarantee does not apply. The response and diagnostics are \OwnedTBD{}.\\
\textbf{Outcome / postconditions} & Acquisition is active without a separate start command.\\
\textbf{Mapped test} & \TestCaseRef{TC-014} (planned).\\
\textbf{Diagram / validation} & \hyperref[Figure:UseCaseDiagram]{UML diagram}; \ValidationRef{VAL-006}.\\
\bottomrule
\end{tabularx}

\clearpage
\section{Recover from Hardware Disconnection}
\begin{tabularx}{\linewidth}{L{3.5cm}Y}
\toprule
\textbf{Identifier} & \UseCaseID{UC-004}; Draft for review.\\
\textbf{Mapped story} & \UserStoryRef{US-005}.\\
\textbf{Mapped requirements} & \RequirementRef{REQ-F-011}; \RequirementRef{REQ-F-012}; \RequirementRef{REQ-F-013}; \RequirementRef{REQ-Q-002}; \RequirementRef{REQ-D-002}.\\
\textbf{Goal} & Resume capture and identify the interruption while preserving host-read bytes.\\
\textbf{Primary actor} & Embedded Engineer. Supporting external actor: downstream processing consumer.\\
\textbf{Preconditions} & Acquisition is active under the \hyperref[Section:ProjectAcquisitionConditions]{agreed acquisition conditions}, and the consumer is able to accept data.\\
\textbf{Trigger} & The FT245R hardware disconnects during acquisition.\\
\textbf{Main flow} & 1. The pipeline remains running and reports the interruption.\par 2. It delivers pending bytes successfully read into host memory before disconnection, unchanged and in order.\par 3. The Embedded Engineer reconnects supported hardware.\par 4. Acquisition resumes automatically with the interruption marked.\\
\textbf{Alternate flow} & If no host-read bytes are pending, the interruption and reconnection steps still apply. If the hardware remains disconnected, acquisition cannot resume.\\
\textbf{Exception flow} & A read failure without confirmed disconnection follows \UseCaseRef{UC-002}. The response to consumer failure and handling of stale device-buffer contents are \OwnedTBD{}.\\
\textbf{Outcome / postconditions} & Acquisition is active after reconnection. The report marks the interruption and makes no claim of complete, uninterrupted acquisition. Device-resident bytes and bytes generated during disconnection are outside the preservation guarantee.\\
\textbf{Mapped test} & \TestCaseRef{TC-015} (planned).\\
\textbf{Diagram / validation} & \hyperref[Figure:UseCaseDiagram]{UML diagram}; \ValidationRef{VAL-007}.\\
\bottomrule
\end{tabularx}

\paragraph{Retired template mappings.}
\label{Section:RetiredTemplateUseCases}
The former submit-request, view-status, and export-history use cases were
removed at the owner's direction. Links in the retained template requirement
and traceability examples target this notice instead of allocating those
unrelated examples to the project use cases. Project requirement and
test allocations remain planned or pending with Joel Murphy and Jason Lefler.
